% (c) Nikita Lisitsa, lisyarus@gmail.com, 2026

\documentclass[10pt]{beamer}

\usepackage[T2A]{fontenc}
\usepackage[russian]{babel}
\usepackage{minted}

\usepackage[most]{tcolorbox}
\tcbuselibrary{minted}

\usepackage{graphicx}
\graphicspath{ {./images/} }

\usetheme{metropolis}

\usepackage{adjustbox}

\usepackage{color}
\usepackage{soul}

\usepackage{hyperref}

\definecolor{codebg}{RGB}{29,35,49}

\setminted{fontsize=\footnotesize}

\setlength\partopsep{-\topsep}

\definecolor{blue}{rgb}{0,0,1}
\definecolor{red}{rgb}{1,0,0}

\usemintedstyle{tango}

\makeatletter
\newcommand{\slideimage}[1]{
  \begin{figure}
    \begin{adjustbox}{width=\textwidth, totalheight=\textheight-2\baselineskip-2\baselineskip,keepaspectratio}
      \includegraphics{#1}
    \end{adjustbox}
  \end{figure}
}
\makeatother

\title{Компьютерная графика}
\subtitle{Практика 10: Normal mapping, material mapping, environment mapping}
\date{2026}

\setbeamertemplate{footline}[frame number]

\begin{document}

\frame{\titlepage}

\begin{frame}[fragile]
\frametitle{Практика 10}
\slideimage{0.png}
\end{frame}

\begin{frame}[fragile]
\frametitle{Задание 1}
Готовимся к normal mapping
\begin{itemize}
\item Выводим нормаль к поверхности вместо цвета:
\usemintedstyle{lightbulb}
\begin{tcblisting}{listing engine=minted, minted language=wgsl, minted options={fontsize=\scriptsize}, listing only, colback=codebg, colframe=codebg, rounded corners}
// let color = ...;
let color = normal * 0.5 + vec3(0.5);
\end{tcblisting}
\end{itemize}
\end{frame}

\begin{frame}[fragile]
\frametitle{Задание 1}
\slideimage{1.png}
\end{frame}

\begin{frame}[fragile]
\frametitle{Задание 2}
Читаем normal map
\begin{itemize}
\item Загружаем текстуру \mintinline{cpp}|normal.jpg| аналогично \mintinline{cpp}|albedo.jpg|
\item Добавляем её в bind group layout и в bind group
\item Добавляем её binding в шейдер
\item Во фрагментном шейдере выводим её значение в качестве цвета
\end{itemize}
\end{frame}

\begin{frame}[fragile]
\frametitle{Задание 2}
\slideimage{2.png}
\end{frame}

\begin{frame}[fragile]
\frametitle{Задание 3}
Вычисляем настоящие нормали
\begin{itemize}
\item Во фрагментном шейдере вычисляем bitangent вектор как \mintinline{wgsl}|bitangent = cross(normal, tangent)|
\item Вычисляем матрицу преобразования из normal map в мировые координаты как \mintinline{wgsl}|tbn = mat3x3f(tangent, bitangent, normal)|
\item Прочитанное из normal map значение переводим из [0, 1] в [-1, 1] и применяем матрицу: \mintinline{wgsl}|real_normal = tbn * (textureSample(...).xyz * 2.0 - vec3f(1.0))|
\item В качестве цвета используем новую нормаль: \mintinline{wgsl}|color = real_normal * 0.5 + vec3(0.5)|
\end{itemize}
\end{frame}

\begin{frame}[fragile]
\frametitle{Задание 3}
\slideimage{3.png}
\end{frame}

\begin{frame}[fragile]
\frametitle{Задание 4}
Используем настоящие нормали для освещения
\begin{itemize}
\item Во фрагментном шейдере возвращаем старое вычисление переменной \mintinline{wgsl}|color|
\item В вычислении освещения используем \mintinline{wgsl}|real_normal| вместо \mintinline{wgsl}|normal|
\end{itemize}
\end{frame}

\begin{frame}[fragile]
\frametitle{Задание 4}
\slideimage{4.png}
\end{frame}

\begin{frame}[fragile]
\frametitle{Задание 5}
Готовимся к reflection + material mapping
\begin{itemize}
\item Загружаем текстуру \mintinline{cpp}|metalness.jpg| аналогично предыдущим (bind group layout, bind group, binding в шейдере)
\item Читаем её в шейдере \mintinline{wgsl}|metalness = textureSample(...)|
\item В качестве цвета выводим \mintinline{wgsl}|vec3f(metalness)|
\end{itemize}
\end{frame}

\begin{frame}[fragile]
\frametitle{Задание 5}
\slideimage{5.png}
\end{frame}

\begin{frame}[fragile]
\frametitle{Задание 6}
Реализуем reflection mapping
\begin{itemize}
\item Загружаем текстуру \mintinline{cpp}|environment_map.jpg| аналогично предыдущим (bind group layout, bind group, binding в шейдере)
\item Во фрагментном шейдере по положению камеры \mintinline{wgsl}|camera.position| и положению пикселя \mintinline{wgsl}|in.world_position| вычисляем направление на камеру
\item С помощью новой нормали вычисляем направление отражённого луча: \mintinline{cpp}|dir = reflect(...)|
\end{itemize}
\end{frame}

\begin{frame}[fragile]
\frametitle{Задание 6}
Реализуем reflection mapping
\begin{itemize}
\item Вычисляем текстурные координаты для environment map:
\usemintedstyle{lightbulb}
\begin{tcblisting}{listing engine=minted, minted language=wgsl, minted options={fontsize=\scriptsize}, listing only, colback=codebg, colframe=codebg, rounded corners}
x = atan2(dir.z, dir.x) / PI * 0.5 + 0.5
y = - atan2(dir.y, length(dir.xz)) / PI + 0.5
\end{tcblisting}
\usemintedstyle{tango}
\item Читаем по этим координатам цвет из environment map: \mintinline{wgsl}|reflected_color = textureSample(...)|
\item В качестве итогого цвета берём \mintinline{wgsl}|mix(color, reflected_color, metalness)| (где \mintinline{wgsl}|color| -- цвет, уже посчитанный из альбедо и освещения)
\end{itemize}
\end{frame}

\begin{frame}[fragile]
\frametitle{Задание 6}
\slideimage{6.png}
\end{frame}

\begin{frame}[fragile]
\frametitle{Задание 7*}
Реализуем environment mapping
\begin{itemize}
\item Создаём новые pipeline и шейдер, которые рисуют один полноэкранный прямоугольник (напр. как в первой практике)
\item Используя матрицу, обратную к view-projection, по ndc-координатам можно восстановить направление в пространстве (см. дополнение к 4ой лекции)
\item Используем это направление, чтобы прочитать значение из environment map и вывести на экран
\item \textbf{NB:} Этот pipeline не должен писать в буфер глубины! Для этого нужен или отдельный render pass, или правильная настройка depth-stencil state
\end{itemize}
\end{frame}

\begin{frame}[fragile]
\frametitle{Задание 7*}
\slideimage{7.png}
\end{frame}

\end{document}
